\section{Introduction}
\label{sec:introduction}

LLM agents increasingly divide work among specialized collaborators, passing plans, intermediate results, and execution requests between roles~\cite{wu2023autogen,hong2024metagpt}. For tasks that change external state, successful collaboration requires preserving the user's authorization throughout this process. Consider a retail customer who approves returning a tea kettle while asking the assistant to hold off on another transaction. A coordinator can delegate both operations, and a specialist can execute the withheld transaction even when the customer conversation remains available. Completing the approved return does not cancel the unauthorized action: a system can accomplish more useful work while also exceeding its permission. Our completed retail study provides a concrete motivation for examining this distinction. Across 150 matched cases for each of two models, individual agents executed an unauthorized transaction in 2 of 300 episodes, compared with 24 of 300 episodes for teams with chat. Approved-task completion was comparable or higher for teams. This result establishes an authorization risk in the evaluated delegation setting and raises a prospective question: could measurements of the individual agents have anticipated that risk?

Existing research provides several components for studying this question. Interactive benchmarks evaluate agents on tool use, task completion, and policy compliance~\cite{liu2024agentbench,yao2024tau}. Analyses of multi-agent trajectories characterize failures in system design, inter-agent alignment, and verification~\cite{cemri2026multi}. Runtime safeguards check or constrain consequential actions during execution~\cite{wang2026agentspec,cuadron2025saber}. PreCast asks how to connect individual measurements to a forecast made before the target team interacts. Ordinary solo evaluation provides limited evidence about responses to delegation: an agent that respects a customer restriction while working alone may behave differently when a teammate assigns the restricted action. Retrospective traces reveal what happened after a team was assembled, while runtime safeguards act once a workflow is underway. A prospective forecast would support an earlier decision: which agents should work together, in which roles, and under which communication protocol?

We formulate \emph{prospective prediction of team authorization risk from controlled individual probes}. Given candidate agents, their assigned roles, a communication protocol, and a task distribution, the goal is to predict the probability that a team executes at least one unauthorized action within a fixed interaction horizon. Individual probes expose one agent at a time to controlled task states and teammate messages, then record its executed actions. For example, a probe can measure whether a specialist obeys a teammate's request when the customer's visible approval covers only the tea-kettle return. Other probes can vary whether a teammate accurately reports permission, claims approval without support, or has already violated a restriction. These measurements characterize responses to potential collaboration contexts. The target evaluation then tests forecasts on previously unobserved team configurations and held-out case groups, with predictions fixed before the corresponding team runs. The central research question is: \emph{can these individual response measurements predict team authorization failures beyond ordinary solo performance and knowledge of the team's model identities?}

Three challenges determine whether this prediction is useful. First, authorization behavior depends on context: a single average violation rate can hide sensitivity to role assignments, peer claims, or conflicting evidence. Second, interaction changes the distribution of those contexts. A team's messages and earlier actions determine what each member encounters, so multiplying isolated failure probabilities cannot generally capture the resulting dependencies. Third, low predicted risk is useful only alongside useful task performance. A forecast that favors agents that refuse every transaction would provide little guidance for deployment. These challenges call for measurements that distinguish behavioral conditions, an explicit account of how individual responses relate to team interactions, and evaluation that tracks both authorization and approved-task completion.

We propose \textsc{PreCast}, a framework for measuring individual responses, forecasting team risk, and using those forecasts to inform team design. Its proposed protocol has three components. \emph{Controlled behavioral elicitation} measures each model--role pair under matched compliant and conflicting peer contexts, including whether a prior peer violation changes the agent's next action. \emph{Team-risk forecasting} combines these response profiles with the specified roles and interaction protocol; development-team data, when used to fit this mapping, are separated from the target evaluation teams. In the retail example, this component must connect a specialist's response to unsupported delegation with the likelihood of encountering such delegation during collaboration. \emph{Prospective evaluation and selection} freezes forecasts, measures calibration and team-risk ranking against simple baselines, and tests whether forecast-guided choices preserve useful completion while reducing violations. In the primary team evaluation, agents generate their own interactions; the controlled peer messages used in probes are diagnostic inputs. The proposed evaluation retains cases with no observed violations and separates executed violations from rejected attempts and interface failures.

This paper develops the following contributions:
\begin{enumerate}
    \item \textbf{A prospective prediction problem.} We define team authorization risk over a fixed interaction horizon and specify the separation between individual elicitation, development-team observations, and held-out team evaluation (Section~\ref{sec:problem}).
    \item \textbf{A behavioral measurement and evaluation protocol.} We propose controlled, role-specific probes and a prospective test of whether their response profiles improve risk forecasts and team selection beyond solo-performance and model-identity baselines (Section~\ref{sec:precast}).
    \item \textbf{Empirical motivation from delegated execution.} Our matched retail study identifies an increase in unauthorized transactions from 2/300 solo episodes to 24/300 team episodes across two models, motivating an evaluation of whether individual probes can anticipate the observed collaboration risk (Section~\ref{sec:experiments}).
\end{enumerate}
The completed study establishes the motivating effect; predictive accuracy and the utility of forecast-guided selection remain the questions for the prospective evaluation.

